% Hafta 14 -- alıştırma gövdesi (ders notu ve ayrı alıştırma kağıdı tarafından \input edilir)
\section{Çözümlü Alıştırmalar}
Her soruyu önce kendiniz deneyin, sonra çözümü okuyun.

\begin{alistirma}
Bir sayının \(\%25\)'inin \(\%40\)'ı \(18\) ise bu sayı kaçtır?
\end{alistirma}
\begin{cozum}
Sayı \(x\) olsun. \(\%25\)'i \(0{,}25x\); bunun \(\%40\)'ı \(0{,}40\cdot 0{,}25x=0{,}10x\).
\[0{,}10x=18\ \Rightarrow\ x=180\]
Kontrol: \(180\)'in \(\%25\)'i \(45\), \(45\)'in \(\%40\)'ı \(18\) \(\checkmark\).
\end{cozum}

\begin{alistirma}
Bir ürünün fiyatına önce \(\%25\) zam, ardından yeni fiyat üzerinden \(\%20\) zam yapılıyor. Son fiyat ilk fiyata göre yüzde kaç artmıştır?
\end{alistirma}
\begin{cozum}
İlk fiyat \(100\) olsun: \(100\xrightarrow{+\%25}125\xrightarrow{+\%20\ (25)}150\).\\
Son fiyat \(150\) \(\Rightarrow\) \(\%50\) artış. (\(\%25+\%20=\%45\) değil; çarpanlar: \(1{,}25\cdot 1{,}2=1{,}5\).)
\end{cozum}

\begin{alistirma}
Bir sınıftaki öğrencilerin \(\%60\)'ı kızdır. Kızların \(\%30\)'u, erkeklerin \(\%45\)'i basketbol oynamaktadır. Sınıfın yüzde kaçı basketbol oynamaktadır?
\end{alistirma}
\begin{cozum}
Sınıf \(100\) kişi olsun: \(60\) kız, \(40\) erkek.
\[60\cdot\frac{30}{100}+40\cdot\frac{45}{100}=18+18=36\ \Rightarrow\ \%36\]
\end{cozum}

\begin{alistirma}
\(\%25\) kârla \(400\) TL'ye satılan bir ürün \(\%10\) zararla satılsaydı kaç TL'ye satılırdı?
\end{alistirma}
\begin{cozum}
Maliyet \(100x\) \(\Rightarrow\) \(\%25\) kârla satış \(125x=400\Rightarrow x=3{,}2\). Maliyet \(=320\) TL.\\
\(\%10\) zararla satış: \(320\cdot 0{,}90=288\) TL.
\end{cozum}

\begin{alistirma}
Bir mağaza etiket fiyatı üzerinden \(\%30\) indirim yaptığında bile maliyet üzerinden \(\%12\) kâr etmektedir. Mağaza indirimden önce maliyet üzerinden yüzde kaç kârla satış yapıyordu?
\end{alistirma}
\begin{cozum}
Maliyet \(100\) \(\Rightarrow\) indirimli satış \(112\). Etiket fiyatı \(E\) ise
\[E\cdot 0{,}70=112\ \Rightarrow\ E=160\]
Etiket fiyatı maliyetin \(\%160\)'ı \(\Rightarrow\) indirimden önce \(\%60\) kâr.
\end{cozum}

\begin{alistirma}
Alkol oranı \(\%40\) olan \(30\) L kolonyaya kaç L saf su eklenirse alkol oranı \(\%25\) olur?
\end{alistirma}
\begin{cozum}
Alkol miktarı değişmez: \(30\cdot 0{,}40=12\) L. Eklenen su \(x\) L:
\[\frac{12}{30+x}=\frac{25}{100}\ \Rightarrow\ 30+x=48\ \Rightarrow\ x=18\ \text{L}\]
\end{cozum}

\begin{alistirma}
Tuz oranı \(\%20\) olan \(40\) kg tuzlu su ile tuz oranı \(\%30\) olan \(60\) kg tuzlu su karıştırılıyor. Karışımdan \(35\) kg su buharlaştırılırsa son tuz oranı yüzde kaç olur?
\end{alistirma}
\begin{cozum}
Tuz: \(40\cdot 0{,}20+60\cdot 0{,}30=8+18=26\) kg; karışım \(100\) kg.\\
Buharlaştırmada tuz değişmez; karışım \(100-35=65\) kg.
\[\frac{26}{65}=\frac{2}{5}=\frac{40}{100}\ \Rightarrow\ \%40\]
\end{cozum}

\begin{alistirma}
\(8000\) TL'nin bir kısmı yıllık \(\%30\), kalanı yıllık \(\%20\) basit faizle bankaya yatırılıyor. \(1\) yıl sonunda toplam \(2000\) TL faiz alındığına göre \(\%30\) faizle kaç TL yatırılmıştır?
\end{alistirma}
\begin{cozum}
\(\%30\)'a yatırılan \(x\) TL, \(\%20\)'ye yatırılan \(8000-x\) TL:
\[\frac{30x}{100}+\frac{20(8000-x)}{100}=2000\ \Rightarrow\ 30x+160000-20x=200000\ \Rightarrow\ x=4000\ \text{TL}\]
Kontrol: \(1200+800=2000\) \(\checkmark\).
\end{cozum}

\kendinizoncesi
\section{Alıştırmalar (Kendiniz Çözün)}

\begin{enumerate}[label=\textbf{\arabic*.},leftmargin=*,itemsep=6pt]
\item \(640\) sayısının \(\%15\)'i kaçtır?
\item \(36\) sayısı \(240\)'ın yüzde kaçıdır?
\item Bir sayının \(\%35\)'i \(84\) ise bu sayı kaçtır?
\item Fiyatı \(1200\) TL olan bir ürüne önce \(\%15\) zam, sonra yeni fiyat üzerinden \(\%10\) indirim yapılıyor. Son fiyat kaç TL'dir?
\item \(\%30\) indirimle \(560\) TL'ye satılan bir ürünün indirimsiz fiyatı kaç TL'dir?
\item Bir dikdörtgenin uzun kenarı \(\%25\) artırılıyor. Alanın değişmemesi için kısa kenar yüzde kaç azaltılmalıdır?
\item Maliyeti \(250\) TL olan bir ürün \(310\) TL'ye satılıyor. Kâr yüzdesi kaçtır?
\item \(5\) tanesi \(20\) TL'ye alınan yumurtaların \(4\) tanesi \(20\) TL'ye satılıyor. Yüzde kaç kâr edilir?
\item Etiket fiyatı \(900\) TL olan bir ürün \(\%20\) indirimle satıldığında maliyet üzerinden \(\%20\) kâr ediliyor. Ürünün maliyeti kaç TL'dir?
\item Şeker oranı \(\%40\) olan \(25\) kg şekerli suya kaç kg su eklenirse şeker oranı \(\%25\) olur?
\item Tuz oranı \(\%20\) olan \(60\) kg tuzlu sudan kaç kg su buharlaştırılırsa tuz oranı \(\%30\) olur?
\item Alkol oranı \(\%15\) olan \(20\) L ile \(\%35\) olan \(30\) L kolonya karıştırılıyor. Yeni alkol oranı yüzde kaçtır?
\item \(22\) ayar \(36\) g altına kaç g saf (\(24\) ayar) altın eklenirse \(23\) ayar altın elde edilir?
\item \(4500\) TL yıllık \(\%40\) basit faizle \(8\) ayda kaç TL faiz getirir?
\item \(6000\) TL basit faizle \(3\) yıl sonunda \(9600\) TL'ye ulaşıyor. Yıllık faiz oranı yüzde kaçtır?
\item \(5000\) TL yıllık \(\%10\) bileşik faizle \(2\) yıl sonunda kaç TL olur?
\end{enumerate}

\subsection*{Cevap Anahtarı}
{\small
\begin{multicols}{2}
\begin{enumerate}[label=\textbf{\arabic*.},leftmargin=*,itemsep=1pt]
\item \(640\cdot 0{,}15=96\).
\item \(\frac{36}{240}\cdot 100=\%15\).
\item \(0{,}35x=84\Rightarrow x=240\).
\item \(1200\cdot 1{,}15\cdot 0{,}9=1242\) TL.
\item \(0{,}7x=560\Rightarrow x=800\) TL.
\item \(1{,}25\cdot k=1\Rightarrow k=0{,}8\): \(\%20\) azaltılmalı.
\item Kâr \(60\); \(\frac{60}{250}\cdot 100=\%24\).
\item \(20\) yumurta: maliyet \(80\), satış \(100\) \(\Rightarrow\) \(\%25\).
\item Satış \(720\); \(1{,}2M=720\Rightarrow M=600\) TL.
\item Şeker \(10\) kg; \(\frac{10}{25+x}=0{,}25\Rightarrow x=15\) kg.
\item Tuz \(12\) kg; \(\frac{12}{60-x}=0{,}3\Rightarrow x=20\) kg.
\item \(\frac{3+10{,}5}{50}=\frac{27}{100}\Rightarrow\%27\).
\item Saf: \(33+x=\frac{23}{24}(36+x)\Rightarrow x=36\) g.
\item \(\frac{4500\cdot 40\cdot 8}{1200}=1200\) TL.
\item \(F=3600=\frac{6000\cdot n\cdot 3}{100}\Rightarrow n=20\): \(\%20\).
\item \(5000\cdot 1{,}1^2=6050\) TL.
\end{enumerate}
\end{multicols}}

% =====================================================
\newcounter{karma}
\newenvironment{karma}[1]{\refstepcounter{karma}\par\medskip\needspace{9\baselineskip}\noindent\textbf{\color{anamavi}Soru \thekarma.}\ {\small\textit{(#1)}}\ }{\par}
\needspace{15\baselineskip}
\section{Finale Hazırlık: Karma Tekrar (Hafta 1--14)}
Bu bölüm final sınavı kapsamındaki tüm konulardan sorular içerir. İlk \(6\) soru tam çözümlüdür; kalan \(10\) soru için bölüm sonunda cevap anahtarı vardır.

\subsection*{Tam Çözümlü Sorular}

\begin{karma}{Hafta 1: faktöriyel}
\(\dfrac{8!-7!}{6!}\) işleminin sonucu kaçtır?
\end{karma}
\begin{cozum}
\(8!=8\cdot 7!\) olduğundan \(8!-7!=7!\,(8-1)=7\cdot 7!\). \ \(7!=7\cdot 6!\):
\[\frac{7\cdot 7\cdot 6!}{6!}=49\]
\end{cozum}

\begin{karma}{Hafta 3: bölünebilme}
\(a\) ve \(b\) rakam olmak üzere, dört basamaklı \(\ab{4a7b}\) sayısı \(45\) ile tam bölünebilmektedir. \(a\)'nın alabileceği değerlerin toplamı kaçtır?
\end{karma}
\begin{cozum}
\(45=5\cdot 9\) ve \(5\) ile \(9\) aralarında asal: sayı hem \(5\)'e hem \(9\)'a bölünmeli.\\
\(5\)'e bölünebilme: \(b=0\) veya \(b=5\). \quad \(9\)'a bölünebilme: \(4+a+7+b=11+a+b\), \(9\)'un katı.\\
\(b=0\): \(11\le 11+a\le 20\) aralığında \(9\)'un tek katı \(18\) \(\Rightarrow a=7\) \ (\(4770=45\cdot 106\)).\\
\(b=5\): \(16\le 16+a\le 25\) aralığında \(9\)'un tek katı \(18\) \(\Rightarrow a=2\) \ (\(4275=45\cdot 95\)).\\
\(a\)'nın değerleri \(7\) ve \(2\) \(\Rightarrow\) toplam \(9\).
\end{cozum}

\begin{karma}{Hafta 7: mutlak değer}
\(|x-2|<5\) eşitsizliğini sağlayan kaç tam sayı vardır? Bu tam sayıların toplamı kaçtır?
\end{karma}
\begin{cozum}
\(|x-2|<5\iff -5<x-2<5\iff -3<x<7\).\\
Tam sayılar: \(-2,-1,0,1,2,3,4,5,6\) \(\Rightarrow\) \(9\) tane. \ Toplam: \((-2)+(-1)+0+1+2+\dots+6=18\).
\end{cozum}

\begin{karma}{Hafta 11: çarpanlara ayırma}
\(x\neq 2\) ve \(x\neq -3\) olmak üzere \(\dfrac{x^2-9}{x^2+x-6}\) ifadesinin en sade biçimi nedir?
\end{karma}
\begin{cozum}
Pay: iki kare farkı, \(x^2-9=(x-3)(x+3)\). \ Payda: çarpımı \(-6\), toplamı \(1\) olan sayılar \(3\) ve \(-2\): \(x^2+x-6=(x+3)(x-2)\).
\[\frac{(x-3)(x+3)}{(x+3)(x-2)}=\frac{x-3}{x-2}\]
\end{cozum}

\begin{karma}{Hafta 14: kâr -- zarar ve indirim}
Bir mağaza bir ürünün etiket fiyatını maliyetin \(\%50\) fazlası olarak belirliyor. Ürün etiket fiyatı üzerinden \(\%20\) indirimle \(360\) TL'ye satılıyor. Ürünün maliyeti kaç TL'dir ve bu satıştan yüzde kaç kâr edilmiştir?
\end{karma}
\begin{cozum}
Etiket \(E\): \ \(E\cdot 0{,}80=360\Rightarrow E=450\) TL. \quad Maliyet \(M\): \ \(1{,}5M=450\Rightarrow M=300\) TL.\\
Kâr: \(360-300=60\) TL \(\Rightarrow\) \(\dfrac{60}{300}\cdot 100=\%20\) kâr.
\end{cozum}

\begin{karma}{Hafta 13: yaş problemi}
Bir annenin bugünkü yaşı, iki çocuğunun yaşları toplamının \(3\) katıdır. \(6\) yıl sonra annenin yaşı, çocuklarının o günkü yaşları toplamının \(2\) katı olacaktır. Anne bugün kaç yaşındadır?
\end{karma}
\begin{cozum}
Çocukların yaşları toplamı bugün \(s\), annenin yaşı \(3s\). \(6\) yıl sonra anne \(3s+6\); iki çocuğun her biri \(6\) yaş büyüyeceğinden toplam \(s+12\).
\[3s+6=2(s+12)\ \Rightarrow\ 3s+6=2s+24\ \Rightarrow\ s=18\]
Anne bugün \(3\cdot 18=54\) yaşındadır. Kontrol: \(6\) yıl sonra \(60=2\cdot 30\) \(\checkmark\).
\end{cozum}

\subsection*{Cevap Anahtarlı Sorular}

\begin{enumerate}[label=\textbf{\arabic*.},leftmargin=*,itemsep=6pt]\setcounter{enumi}{6}
\item \textit{(Hafta 2)} \(180\) sayısının kaç pozitif böleni vardır?
\item \textit{(Hafta 3)} Kenarları \(48\) m ve \(72\) m olan dikdörtgen bir bahçe, eşit büyüklükte kare parsellere ayrılacaktır. En az kaç parsel elde edilir?
\item \textit{(Hafta 4)} \(0{,}5+\dfrac{2}{3}\cdot\dfrac{3}{4}\) işleminin sonucu kaçtır?
\item \textit{(Hafta 5)} \(3(x-2)+4=2x+5\) denkleminin çözümü nedir?
\item \textit{(Hafta 6)} Bir otoparkta otomobil ve motosikletlerden oluşan \(40\) araç vardır. Tekerleklerin toplam sayısı \(130\) olduğuna göre kaç motosiklet vardır? (Otomobil \(4\), motosiklet \(2\) tekerleklidir.)
\item \textit{(Hafta 9)} \(9^{x}=27^{x-1}\) ise \(x\) kaçtır?
\item \textit{(Hafta 10)} \(\dfrac{6}{\sqrt{3}}+\sqrt{27}\) işleminin sonucu kaçtır?
\item \textit{(Hafta 12)} \(5\) sayının aritmetik ortalaması \(18\)'dir. Bu sayılardan biri çıkarılınca kalan \(4\) sayının ortalaması \(16\) oluyor. Çıkarılan sayı kaçtır?
\item \textit{(Hafta 14)} Tuz oranı \(\%20\) olan \(30\) kg tuzlu suya \(10\) kg tuz ekleniyor. Yeni tuz oranı yüzde kaçtır?
\item \textit{(Hafta 14)} \(7200\) TL yıllık \(\%25\) basit faizle \(120\) günde kaç TL faiz getirir? (\(1\) yıl \(=360\) gün)
\end{enumerate}

\subsection*{Cevap Anahtarı (Karma Tekrar)}
{\small
\begin{multicols}{2}
\begin{enumerate}[label=\textbf{\arabic*.},leftmargin=*,itemsep=1pt]
\item \(49\) \ (tam çözümlü)
\item \(9\) \ (tam çözümlü)
\item \(9\) tam sayı; toplam \(18\) \ (tam çözümlü)
\item \(\frac{x-3}{x-2}\) \ (tam çözümlü)
\item Maliyet \(300\) TL; \(\%20\) kâr \ (tam çözümlü)
\item \(54\) \ (tam çözümlü)
\item \(180=2^2\cdot 3^2\cdot 5\Rightarrow 3\cdot 3\cdot 2=18\).
\item EBOB\((48,72)=24\); \(\frac{48}{24}\cdot\frac{72}{24}=2\cdot 3=6\) parsel.
\item \(0{,}5+0{,}5=1\).
\item \(3x-2=2x+5\Rightarrow x=7\).
\item \(4a+2m=130,\ a+m=40\Rightarrow a=25,\ m=15\).
\item \(3^{2x}=3^{3x-3}\Rightarrow 2x=3x-3\Rightarrow x=3\).
\item \(2\sqrt{3}+3\sqrt{3}=5\sqrt{3}\).
\item \(5\cdot 18-4\cdot 16=90-64=26\).
\item \(\frac{6+10}{30+10}=\frac{16}{40}\Rightarrow\%40\).
\item \(\frac{7200\cdot 25\cdot 120}{36000}=600\) TL.
\end{enumerate}
\end{multicols}}
